% ----------------------------------------------------------
\chapter{Metodologia}
% ----------------------------------------------------------

\section[Materiais]{Materiais}

Nesta seção, são listados os equipamentos de \textit{hardware} e as ferramentas, linguagens, bibliotecas e plataformas de \textit{software} empregados na concepção, no desenvolvimento e na validação da aplicação, de modo a permitir a reprodução do ambiente de trabalho. As decisões de arquitetura que motivaram a escolha dos itens mais relevantes são discutidas e justificadas na seção de Métodos, a seguir.

\subsection[Equipamentos]{Equipamentos}

\begin{itemize}
    \item Notebook: processador Intel Core i5, 1.6 GHz, 8 GB de memória RAM e SSD de 240 GB, com \textit{dual boot} para os sistemas operacionais Windows 10 e Linux Mint.
    \item PC: processador AMD Ryzen 5 5600G, 3.9 GHz, 32 GB de memória RAM e SSD de 477 GB, com sistema operacional Windows 11.
\end{itemize}

\subsection[Análise, Modelagem e Projeto]{Análise, Modelagem e Projeto}

\begin{itemize}
    \item Unified Modeling Language (UML) \cite{UML2015} --- elaboração dos diagramas de caso de uso do sistema.
    \item Excalidraw, com a biblioteca \textit{Shapes for UML and ER Diagrams} --- criação dos diagramas de arquitetura e dos modelos de dados.
\end{itemize}

\subsection[Persistência de dados]{Persistência de dados}

\begin{itemize}
    \item PostgreSQL \cite{PostgreSQL2024} --- persistência dos dados relacionais (usuários, projetos e permissões de acesso).
    \item MongoDB \cite{MongoDB2021} --- persistência da estrutura dos diagramas ER.
    \item DBeaver --- administração, consulta e manipulação dos dados no PostgreSQL e no MongoDB.
\end{itemize}

\subsection[Implementação]{Implementação}

\begin{itemize}
    \item Git \cite{Git2025} --- controle de versão do código-fonte.
    \item GitHub \cite{GitHub2024} --- hospedagem e repositório central do código-fonte.
    \item GitHub Copilot --- apoio à revisão de código em parte dos \textit{pull requests}.
    \item Docker \cite{Docker2024} --- containerização do ambiente de desenvolvimento e dos bancos de dados.
    \item IntelliJ IDEA \cite{IntelliJ2024} --- IDE para o desenvolvimento do \textit{back-end} em Java e Spring.
    \item WebStorm \cite{WebStorm2024} --- IDE para o desenvolvimento do \textit{front-end} em Angular.
    \item Java \cite{Java2024} --- linguagem de programação do \textit{back-end}.
    \item Spring Boot \cite{SpringBoot2024} --- \textit{framework} de estruturação do \textit{back-end} e das APIs REST.
    \item Spring WebSocket \cite{WebSockets2024} --- comunicação em tempo real entre \textit{back-end} e clientes.
    \item Spring Security \cite{SpringSecurity2024} --- autenticação e autorização da aplicação, em conjunto com JWT.
    \item Spring Data JPA \cite{SpringDataJPA2024} --- camada de acesso a dados do PostgreSQL.
    \item Spring Data MongoDB \cite{SpringDataMongoDB2024} --- camada de acesso a dados do MongoDB.
    \item Angular \cite{Angular2022} --- construção da interface do usuário como \textit{Single-Page Application}.
    \item GoJS \cite{GoJS2024} --- renderização e manipulação interativa dos diagramas ER.
    \item stompjs e SockJS \cite{StompJS2023} --- comunicação via \textit{WebSocket} no lado do cliente.
    \item Node Package Manager (npm) \cite{NPM2024} --- gestão das dependências do \textit{front-end}.
\end{itemize}

\subsection[Testes]{Testes}
\label{subsec:testes}

\begin{itemize}
    \item JUnit e Mockito \cite{JUnit2023} \cite{Mockito} --- testes unitários da lógica de negócio no \textit{back-end}.
    \item Jasmine e Karma \cite{Jasmine2024} \cite{Karma2024} --- testes unitários dos componentes e serviços do \textit{front-end}.
    \item JaCoCo \cite{JaCoCo2024} --- medição da cobertura de código dos testes unitários do \textit{back-end}.
    \item SonarQube \cite{SonarQube2024} --- análise estática de qualidade e cobertura de código do \textit{back-end}.
    \item Postman \cite{Postman2024} --- testes manuais das APIs REST do \textit{back-end}.
    \item Google Chrome e Mozilla Firefox \cite{GoogleChrome} \cite{MozillaFirefox} --- testes manuais de compatibilidade da interface do usuário.
\end{itemize}

\section[Métodos]{Métodos}
Nesta seção, são apresentados os métodos e as abordagens que nortearam o planejamento, a análise, o projeto e a implementação da aplicação. A condução do trabalho foi baseada em uma combinação de práticas conhecidas da Engenharia de \textit{Software} e da pesquisa científica, adaptadas para atender suas especificidades.

\subsection[Abordagem Metodológica]{Abordagem Metodológica}
Para a condução deste projeto, adotou-se uma abordagem metodológica híbrida, que integra conceitos do método científico, elementos do modelo em espiral de Boehm e práticas da metodologia Scrum.

\subsubsection[Método Científico]{Método Científico}
O método científico serviu como alicerce estrutural do trabalho, definindo as etapas macro de observação de um problema (a necessidade de uma ferramenta de modelagem ER colaborativa em tempo real), formulação de uma hipótese (a viabilidade de construir tal aplicação com tecnologias \textit{web}), experimentação (o desenvolvimento da solução) e análise dos resultados.

\subsubsection[Modelo em Espiral de Boehm]{Modelo em Espiral de Boehm}
O modelo foi empregado conceitualmente para guiar o ciclo de vida do projeto de forma iterativa e incremental. Cada ciclo envolveu planejamento, análise de riscos (como a complexidade da sincronização em tempo real), engenharia (codificação e testes) e avaliação pelo desenvolvedor, permitindo o refinamento contínuo da aplicação a cada iteração.

\subsubsection[Metodologia Scrum]{Metodologia Scrum}
A metodologia foi utilizada de forma adaptada para o gerenciamento das atividades de desenvolvimento. O trabalho foi decomposto em épicos, representando as grandes funcionalidades, que por sua vez foram detalhados em tarefas menores e mais específicas, facilitando o acompanhamento do progresso e a organização do trabalho em ciclos curtos de desenvolvimento.

\subsection[Análise e Projeto do Sistema]{Análise e Projeto do Sistema}
A fase de análise e projeto foi fundamental para traduzir os requisitos do sistema em um modelo funcional e técnico. Para a definição dos requisitos funcionais e do escopo do sistema, foram elaborados diagramas de Caso de Uso. Essa técnica da \textit{Unified Modeling Language} (UML) permitiu identificar os atores e suas interações com o sistema, delimitando claramente as funcionalidades a serem implementadas, como criar um projeto, convidar colaboradores, adicionar uma entidade ao diagrama, entre outras. A partir dessa análise, foram tomadas decisões arquiteturais cruciais.

\subsubsection[Arquitetura Cliente-Servidor]{Arquitetura Cliente-Servidor}
%
% Explicar porque escolheu a arquitetura Cliente-Servidor
%
Optou-se pela arquitetura macro do tipo cliente-servidor, padrão para aplicações \textit{web} modernas. Nessa arquitetura, o \textit{front-end} (cliente), executado no navegador do usuário, é responsável pela interface e interação, enquanto o \textit{back-end} (servidor) centraliza a lógica de negócio, o controle de acesso e a persistência dos dados.

\subsubsection[Protocolo de Comunicação em Tempo Real]{Protocolo de Comunicação em Tempo Real}

Para atender ao requisito central de colaboração em tempo real, o protocolo WebSocket foi escolhido para a comunicação contínua entre cliente e servidor. A comunicação síncrona do protocolo HTTP, baseada em um ciclo de requisição e resposta estritamente iniciado pelo cliente, mostra-se ineficiente para este cenário. Soluções paliativas que utilizam HTTP, como o \textit{polling} (onde o cliente consulta o servidor em intervalos regulares), geram uma alta latência e um tráfego de rede excessivo, pois cada consulta reinicia uma conexão e carrega metadados que são, na maioria das vezes, desnecessários.

O WebSocket, em contrapartida, estabelece uma conexão única, persistente e bidirecional  após um \textit{handshake} inicial. Essa arquitetura permite que o servidor envie atualizações aos clientes de forma proativa e instantânea (\textit{server push}), com baixa sobrecarga de dados. Essa característica é fundamental para a viabilidade de sistemas colaborativos, onde a propagação de alterações deve ter a menor latência possível para garantir uma experiência de uso fluida e consistente entre os múltiplos usuários. Conforme \citeonline{Souza2023}, em aplicações que demandam alta interatividade e troca constante de pequenos pacotes de dados, o WebSocket apresenta um desempenho superior e um custo computacional significativamente menor quando comparado às abordagens baseadas em HTTP.

\subsection[Modelagem das Bases de Dados]{Modelagem das Bases de Dados}
A estratégia de persistência de dados adotada neste projeto foi a abordagem híbrida, também conhecida como persistência poliglota (\textit{polyglot persistence}). Essa decisão consistiu em combinar um sistema de gerenciamento de banco de dados relacional e um não relacional, a fim de utilizar a tecnologia mais adequada para cada tipo de dado manipulado pela aplicação, melhorando tanto o desempenho quanto a integridade da informação.

\subsubsection[Dados Relacionais]{Dados Relacionais}
%
% Explicar porque escolheu o PostgreSQL
%
Para o armazenamento de dados com estrutura bem definida e que exigem integridade referencial, como as informações de usuários, projetos e credenciais de acesso, optou-se pelo SGBD PostgreSQL. Sua conformidade com as propriedades ACID (Atomicidade, Consistência, Isolamento e Durabilidade) e seu sistema de controle de transações foram fatores determinantes para garantir a consistência e a segurança dos dados do sistema. A modelagem dessas informações foi formalizada por meio de um diagrama Entidade-Relacionamento (ER).

\subsubsection[Dados Não Relacionais]{Dados Não Relacionais}
%
% Falar que o Spring tem um módulo do MongoDB, o que facilitaria no desenvolvimento
%
Em contrapartida, para a persistência dos dados que compõem os diagramas ER do protótipo (como entidades, atributos, relacionamentos e suas respectivas posições na área de desenho) a escolha recaiu sobre o MongoDB. A natureza colaborativa e em tempo real da aplicação impõe um requisito de destaque para operações de leitura e escrita concorrentes, onde múltiplos usuários podem consultar e alterar o mesmo diagrama. Esta decisão é respaldada por estudos que demonstram a superioridade do MongoDB em relação ao PostgreSQL \cite{Politowski2014}. O modelo de documentos permite que a estrutura de um diagrama seja armazenada de forma coesa e hierárquica, otimizando as consultas e simplificando a lógica de sincronização necessária para a colaboração em tempo real.

\subsection[Dialeto SQL]{Dialeto SQL}

Uma das funcionalidades centrais do protótipo é a capacidade de exportar e importar o diagrama ER como um \textit{script} SQL. Para viabilizar essas operações, foi necessário definir um dialeto SQL de referência, já que diferentes SGBDs possuem particularidades sintáticas próprias em suas instruções DDL.

O dialeto adotado foi o {MySQL. O MySQL é um sistema de gerenciamento de banco de dados relacional de código aberto, mantido pela Oracle Corporation, e é amplamente reconhecido como um dos SGBDs mais utilizados no mundo para aplicações \textit{web} \cite{OracleMySQL}. Sua sintaxe DDL é amplamente difundida no mercado de desenvolvimento \textit{web}, o que torna o dialeto MySQL uma escolha adequada para maximizar a interoperabilidade dos \textit{scripts} gerados e importados pelo protótipo com outros ambientes e ferramentas do ecossistema de desenvolvimento de \textit{software}.

\subsection[Padrão de Arquitetura de Software e de Código]{Padrão de Arquitetura de Software e de Código}
Com a arquitetura macro definida, foram adotados padrões específicos para a organização interna do código, tanto no \textit{back-end} quanto no \textit{front-end}.

\subsubsection[Model-View-Controller]{Model-View-Controller}
%
% Falar que escolheu o Spring pelo suporte ao WebSocket
% Falar que o Spring usa o MVC como padrão e por este motivo adotou-se a linguagem Java e o Angular, associando a 
% divisão de responsabilidades com o princípio do Clean Code
% 
% https://sabrinabm94.medium.com/usando-o-mvc-no-angular-13081366019a
% https://www.amazon.com.br/Iniciando-com-JavaScript-fundamentos-Aprendendo/dp/B0D25257DD
%
A escolha do padrão de arquitetura \textit{Model-View-Controller} (MVC) foi uma decisão estratégica para a organização do código-fonte, sendo aplicada tanto no \textit{back-end} quanto no \textit{front-end} da aplicação. Este padrão promove a separação de responsabilidades, dividindo a aplicação em três componentes interconectados, o que foi útil para gerenciar a complexidade do sistema.

\begin{itemize}
    \item \textit{Model}: Camada de manipulação dos dados.
    \item \textit{View}: Camada de interação do usuário.
    \item \textit{Controller}: Camada de controle, responsável por ligar o \textit{model} e a \textit{view}, fazendo com que os \textit{models} possam ser repassados para as \textit{views} e vice-versa.
\end{itemize}

A adoção do padrão MVC em ambas as camadas do sistema não foi uma mera formalidade, mas uma escolha deliberada para obter vantagens concretas no processo de desenvolvimento. Essa abordagem está alinhada com as melhores práticas da Engenharia de \textit{Software}, que destacam o MVC como um padrão que promove a coesão e o baixo acoplamento entre as partes de um sistema. Conforme aponta a literatura, a separação de interesses facilita o desenvolvimento paralelo, a reutilização de componentes e, crucialmente, a testabilidade e a manutenção do \textit{software} \cite{EngenhariaSoftware2021}.

\begin{itemize}
    \item \textit{Back-end}: O \textit{framework} Spring implementa nativamente uma variação do MVC. As classes que usam uma determinada anotação atuam como Controladores, recebendo requisições HTTP. A lógica de negócio e os objetos de domínio representam o Modelo. A Visão é materializada pela resposta serializada, tipicamente em formato JSON, enviada ao cliente.
    \item \textit{Front-end}: O \textit{framework} Angular organiza a aplicação de forma análoga. Os componentes são formados por um \textit{template} HTML, uma classe que contém a lógica e o estado, e os serviços que gerenciam os dados da aplicação.
\end{itemize}

\subsubsection[Clean Code]{Clean Code}
A qualidade interna do código-fonte foi uma adotada durante todo o ciclo de desenvolvimento. Para tanto, foram escolhidos de forma sistemática os princípios de Código Limpo, conforme popularizados por \citeonline{CleanCode2009}, com o objetivo de produzir uma base de código legível e manutenível. A aplicação desses princípios se manifestou principalmente nas seguintes práticas:

\begin{itemize}
    \item Nomes de variáveis e métodos significativos
    \item Métodos pequenos e com responsabilidade única
    \item Código autoexplicativo
\end{itemize}

O primeiro princípio foi aplicado promovendo atenção especial à nomenclatura para que o propósito de cada elemento do código fosse explícito. Em vez de nomes genéricos, utilizaram-se identificadores descritivos. Da mesma forma, os métodos foram nomeados para expressar claramente sua ação e seu resultado. Essa prática elimina a ambiguidade e torna o código mais intuitivo.

O segundo princípio foi aplicado na decomposição da lógica de negócio. Métodos e funções foram projetados para executar uma única tarefa e fazê-la bem. Por exemplo, um controlador de API REST que recebe uma requisição para adicionar uma nova entidade ao diagrama não executa a validação, a persistência no banco e a notificação via WebSocket dentro de um único método monolítico. Em vez disso, ele delega cada uma dessas responsabilidades a camadas de serviço distintas: um método para validar os dados, um para orquestrar a lógica de negócio e a persistência, e outro para notificar os clientes conectados. Essa separação torna o sistema mais modular, testável e fácil de depurar.

O terceiro princípio foi majoritariamente empregado na escrita de um código que, por sua própria estrutura e clareza, documenta seu funcionamento, minimizando a necessidade de comentários explicativos. A combinação de nomes significativos e funções com responsabilidade única resultou em um código que se assemelha a uma narrativa legível. Comentários foram utilizados de forma criteriosa, reservados para justificar decisões de projeto não triviais ou para explicar algoritmos complexos, em vez de parafrasear o que uma linha de código já expressa claramente.

\subsection[Estratégia de Controle de Versão]{Estratégia de Controle de Versão}
A gestão do código-fonte é um pilar para a qualidade e a rastreabilidade de qualquer projeto de software. Para este trabalho, foi empregado o Sistema de Controle de Versão (VCS) \textit{git}, com o código hospedado na plataforma GitHub \cite{Git2014}. A estratégia adotada foi um modelo híbrido, que combina a agilidade do \textit{Trunk-Based Development} como fluxo principal, complementado por uma adaptação simplificada do \textit{GitFlow} para o desenvolvimento de funcionalidades específicas.

A abordagem primária do \textit{Trunk-Based Development}, na qual as alterações são integradas diretamente na ramificação principal (master), foi uma escolha deliberada para promover um ritmo de desenvolvimento mais acelerado e simplificar o gerenciamento de versões. A viabilidade dessa abordagem foi sustentada pelo contexto do projeto, que conta com um único desenvolvedor. Essa decisão encontra respaldo no estudo de \citeonline{Lopes2025}, que aponta que a abordagem \textit{Trunk-Based} é particularmente eficaz em projetos de ritmo acelerado, com equipes menores e mais experientes, onde a complexidade de gerenciamento de múltiplas ramificações se torna um \textit{overhead} desnecessário.

Para o desenvolvimento de funcionalidades mais complexas ou para a correção de \textit{bugs} que exigissem um trabalho mais extenso, foi adotada uma adaptação do modelo \textit{GitFlow}. Neste fluxo simplificado, uma ramificação de funcionalidade (\textit{feature branch}) era criada a partir da master. Esta prática permitiu implementar e testar novas capacidades de forma isolada, sem comprometer a estabilidade da ramificação principal. Após a conclusão e validação, a \textit{feature branch} era integrada de volta à master. Esse isolamento facilitou a organização das tarefas e garantiu uma camada adicional de validação, além de aprimorar a rastreabilidade do histórico de desenvolvimento, uma vez que cada conjunto de alterações significativas fica contido em sua própria ramificação \cite{Rivero2021}.

A combinação dessas duas estratégias proporcionou o melhor de ambos os mundos para o escopo deste trabalho: a velocidade e a simplicidade do \textit{Trunk-Based Development} para o progresso diário e a segurança e organização das ramificações de funcionalidade para as implementações mais críticas.

\subsection[Autenticação e Autorização]{Autenticação e Autorização}
A segurança da aplicação foi implementada utilizando uma combinação de tecnologias e estratégias, visando proteger os dados e garantir o acesso adequado aos recursos.  O Spring Security, um \textit{framework} utilizado no ecossistema Java, foi o pilar central da autenticação e autorização \cite{SpringSecurity2024}.

Para a autenticação, foi adotado o padrão JWT.  Após a autenticação bem-sucedida do usuário (verificação de credenciais), o servidor gera um \textit{token} JWT, que é então enviado ao cliente. Este \textit{token} contém informações sobre o usuário autenticado e suas permissões.

Os \textit{tokens} JWT foram armazenados em \textit{cookies} HTTP no cliente. Essa abordagem permite que o navegador gerencie automaticamente o envio do \textit{token} em cada requisição subsequente ao servidor, simplificando a implementação e melhorando a segurança em relação ao armazenamento em \textit{localStorage}, que é mais vulnerável a ataques XSS (\textit{Cross-Site Scripting}) \cite{Bulgakova2020}.

No \textit{back-end}, o Spring Security intercepta todas as requisições, validando a autenticidade do \textit{token} JWT presente no \textit{cookie}.  Caso o \textit{token} seja válido, o Spring Security extrai as informações do usuário e suas permissões, permitindo o acesso aos recursos autorizados.

No \textit{front-end}, Angular Guards atuam para proteger rotas e componentes, impedindo o acesso não autorizado a determinadas áreas da aplicação, de acordo com \citeonline{Oliveira2024}.  Por exemplo, apenas usuários com ua determinada permissão podem acessar a página de gerenciamento de usuários.

Essa abordagem em camadas, combinando Spring Security, JWT, \textit{Cookies} e Angular Guards, oferece uma solução de segurança, protegendo a aplicação contra acesso não autorizado e garantindo a integridade dos dados.

\subsection[Estratégia de Concorrência]{Estratégia de Concorrência}
A funcionalidade de colaboração em tempo real introduz a necessidade de uma estratégia para lidar com a concorrência.  Múltiplos usuários editando o mesmo elemento simultaneamente podem levar a conflitos e perda de dados se não forem devidamente gerenciados.  De acordo com \citeonline{Hipolito2024}, para mitigar esse problema, foi implementado um mecanismo de \textit{locks} (travas) com armazenamento em memória.

A estratégia de \textit{locks} funciona da seguinte forma: antes de um usuário iniciar a edição de uma seção específica de um elemento, um \textit{lock} correspondente a essa seção é adquirido e impede que outros usuários editem a mesma seção simultaneamente.  O \textit{lock} é liberado assim que o usuário termina de editar a seção, permitindo que outros usuários a acessem.

A utilização de armazenamento em memória garante um acesso rápido e eficiente às informações de travamento, minimizando a latência e, consequentemente, o impacto na experiência do usuário durante a colaboração em tempo real.  A implementação específica dos \textit{locks} foi realizada utilizando um recurso interno da linguagem Java, que oferece ferramentas eficientes para gerenciamento de concorrência em memória.

Embora essa abordagem mitigue os conflitos de edição, ela introduz a possibilidade de \textit{deadlocks} (travamentos mútuos). Com a finalidade de evitar isso, foi implementada uma estratégia de \textit{timeout} para os \textit{locks}. Se uma trava não for liberada dentro de um período de tempo determinado, ela é automaticamente liberada pelo sistema, prevenindo que outros usuários fiquem indefinidamente bloqueados \cite{Filippo2005}.

Essa estratégia de concorrência, baseada em travas em memória com \textit{timeouts}, oferece um equilíbrio entre a proteção dos dados e a manutenção de uma experiência de colaboração em tempo real fluida e responsiva.

\subsection[Design da Interface do Usuário]{Design da Interface do Usuário}
O \textit{design} da \textit{interface} da aplicação foi elaborado para garantir uma experiência de usuário intuitiva e eficaz, combinando a familiaridade de ferramentas de modelagem de dados estabelecidas com a usabilidade e a estética de aplicações \textit{web} modernas.

A inspiração inicial para a estrutura geral da interface foi o MySQL Workbench. Esta escolha visou facilitar a transição para usuários já familiarizados com este ambiente de modelagem de dados. No entanto, a \textit{interface} foi adaptada e enriquecida com padrões de \textit{design} modernos comumente encontrados em aplicações \textit{web}, tais como:

\begin{itemize}
    \item Menu Lateral: Um menu lateral fixo oferece acesso rápido e fácil a todas as funcionalidades principais da aplicação.
    \item \textit{Modals}: Janelas \textit{modals} são utilizadas para apresentar informações adicionais ou solicitar confirmação do usuário em tarefas específicas, mantendo o contexto da tela principal.
    \item Elementos Visuais: Utilização de elementos visuais consistentes e modernos para guiar o usuário e melhorar a clareza da informação.
\end{itemize}

Um componente central da \textit{interface} é a visualização e manipulação interativa dos diagramas. Para essa funcionalidade, foi integrada a biblioteca GoJS. Ela permite a criação de diagramas complexos e dinâmicos, com recursos como arrastar e soltar, \textit{zoom}, e personalização visual, proporcionando uma experiência de modelagem de dados interativa. Essa abordagem torna a criação e visualização de diagramas uma parte central da experiência do usuário, facilitando a compreensão e manipulação das estruturas de dados.

A combinação do design conhecido do MySQL Workbench com elementos modernos de aplicações \textit{web} e a biblioteca GoJS indicam que uma \textit{interface} intuitiva e eficaz para a modelagem de dados e a colaboração em tempo real possa ser construída.



\subsection[Idioma do Projeto]{Idioma do Projeto}

A língua inglesa exerce uma posição dominante na área da Tecnologia da Informação. Conforme \cite{Alves2022}, a globalização e a competitividade do mercado tornaram o inglês praticamente nativo no setor tecnológico, sendo a língua predominante nas áreas de divulgação científica, comércio e relações internacionais. Essa realidade se manifesta de forma concreta na prática do desenvolvimento de \textit{software}: os comandos e as palavras reservadas de todas as principais linguagens de programação são estruturados em inglês, assim como os métodos ágeis amplamente adotados pela indústria, como o Scrum e o Kanban.

Diante desse contexto, a adoção do inglês como idioma padrão para todos os artefatos técnicos produzidos neste projeto foi uma decisão deliberada, fundamentada nas boas práticas da Engenharia de \textit{Software}. Esta escolha está em consonância com os princípios do Código Limpo (\textit{Clean Code}), que preconiza o uso de nomes significativos e autoexplicativos \cite{CleanCode2009}. A utilização do inglês garante que a nomenclatura empregada no código-fonte, nos esquemas de banco de dados e nos diagramas seja consistente com a terminologia das tecnologias, bibliotecas e \textit{frameworks} adotados no projeto, evitando ambiguidades e facilitando a compreensão e a manutenção da base de código.

Esta padronização foi aplicada de forma abrangente nos seguintes artefatos do projeto:

\begin{itemize}
    \item Código-fonte: nomes de classes, métodos, variáveis e atributos foram definidos em inglês tanto no \textit{back-end} (Java/Spring) quanto no \textit{front-end} (Angular/TypeScript), mantendo coerência com os próprios nomes das APIs e as convenções dos \textit{frameworks} utilizados.
    \item Banco de dados: os nomes das tabelas, colunas e restrições no PostgreSQL, bem como os nomes das coleções e campos dos documentos no MongoDB, foram definidos em inglês. Essa decisão garante consistência com as anotações do Spring Data JPA e do Spring Data MongoDB, que realizam o mapeamento objeto-relacional e objeto-documento a partir dos identificadores declarados no código-fonte.
    \item Diagramas: os nomes de entidades, atributos e relacionamentos nos diagramas ER foram especificados em inglês, garantindo uniformidade com os demais artefatos técnicos do projeto.
    \item \textit{Endpoints} da API REST: os caminhos dos recursos e as operações expostas pela API foram nomeados em inglês, seguindo as convenções estabelecidas pelo estilo arquitetural REST.
    \item Controle de versão: as mensagens de \textit{commit} e os nomes de ramificações (\textit{branches}) no repositório \textit{git} foram redigidos em inglês, promovendo a rastreabilidade e a clareza do histórico de desenvolvimento.
\end{itemize}

A adoção uniforme do inglês como idioma técnico em todos os artefatos do projeto reduziu o atrito cognitivo entre as camadas da aplicação, facilitou a leitura e a manutenção do código e posicionou o trabalho em conformidade com os padrões internacionais da indústria de desenvolvimento de \textit{software}.
